\section{What is being computed}\label{cx:sec:computed}

Collecting the chapter: given $\zeta$ on the line $\Re s=1$ to an absolute precision on the scale of $u_1$, which for the linear inverse costs $(1-\theta')\log_2T$ bits near height $T$ (Proposition~\ref{cx:prop:precision}), which is the expensive part and the only part that grows faster than the output, the EM iteration and its accelerated forms compute the information projection of the level one Thorin measure onto the cone of Cauchy blurred positive measures at level $\theta'$. When $\zeta$ has no zeros to the right of $\theta'$ the limit is the Thorin measure $U_{\theta'}$ of $\xi(\theta')/\xi(\theta'+\sqrt s)$ (Proposition~\ref{cx:prop:what}(a), under its conditions on the windows). Under the Riemann hypothesis this is the Brownian image of the zero measure stopped at the first passage of level $\theta'-\tfrac12$, and the zeros are its peaks. When there is a zero to the right of $\theta'$, the residual of the grid problems stays bounded away from zero as the grid is refined (Proposition~\ref{cx:prop:what}(b)). None of this reaches the arithmetic, which lives in the evaluation of $\zeta$ and costs polynomial time in $T$ by every known method.

In the rest of this part, Theorems~\ref{cx:thm:certified} and~\ref{cx:thm:globalmin} and the EM runs of Section~\ref{cx:sec:half} were recomputed; the three routes of Section~\ref{cx:sec:level78}, the mass deficit and level mass tables, the tables of $U_{3/4}$ and the timings are the author's, not rerun. The author's figures are not reproduced; the text says what they show.

\section{Thorin measures at level 7/8}\label{cx:sec:level78}

\subsection{Thorin measures, not primes}

What the backward flow computes well is Thorin measures. At every level $\theta\ge1$ the Thorin measure of $\xi(\theta)/\xi(\theta+\sqrt s)$ has the explicit density $\pi^{-1}\Re(\xi'/\xi)(\theta+iv)$, with the prime sum $-\sum\Lambda(n)n^{-\theta}\cos(v\log n)$ inside it (absolutely convergent for $\theta>1$; at $\theta=1$ it stands for the value of $\zeta'/\zeta$ on the line), computable on one line with no zeros and no inversion (Chapter~\ref{chap:lo}, Theorem~\ref{lv:thm:thorin}). Below level one, above the zero free abscissa, the Thorin measure is the EM limit from the level one data, and at $\theta'=\tfrac12$ it is resolved down to the atoms: the thousand zero run of Section~\ref{cx:sec:thousand} computes the level one half Thorin measure of the reciprocal $\xi$ function from the Euler product side. Under the Riemann hypothesis the measure at level $\theta>\tfrac12$ is the Cauchy smoothing of the zero measure at width $\theta-\tfrac12$, so a computed Thorin measure is to be compared with that smoothed density, not with the atomic measure in total variation, where the two are mutually singular whatever the truth.

What the flow does not compute is primes. They are the spectrum of the level one measure, point masses at $\pm\log n$ of weight proportional to $-\Lambda(n)/n$, put in through the Euler product when $u_1$ is evaluated; reading them off (Chapter~\ref{chap:lo}) recovers what was supplied. The reverse passage, from zeros to primes, is the classical explicit formula for $\psi(x)$, which nothing here changes.

\subsection{The level $7/8$ measure}

The quasi Riemann hypothesis (Section~\ref{lv:sec:qrh}) is the zero free half plane $\Re s>7/8$, announced in October 2026 by OpenAI with a formalisation in Lean, and by Liu, not yet refereed \cite{OAI26a,OAI26b,Liu26a}. It makes $\theta=7/8$ the lowest level at which the Thorin representation is currently claimed; every statement below that rests on it is conditional on it, as in Chapter~\ref{chap:lv}:
\begin{equation}\label{cx:eq:F78}
\frac{\xi(7/8)}{\xi(7/8+\sqrt s)}=\exp\Bigl(-\int_0^\infty\log\Bigl(1+\frac su\Bigr)U_{7/8}(du)\Bigr),
\end{equation}
with $U_{7/8}$ a positive Thorin measure (Theorem~\ref{lv:thm:thorin}). We compute it three ways, the agreement being the check. Throughout, $u(v)$ denotes the density of the symmetric measure $\mathfrak U_{7/8}$ in the variable $v$; the Thorin measure $U_{7/8}$ in the variable $u=v^2$ is the image of the restriction of $\mathfrak U_{7/8}$ to $v>0$, so that $\int g(u)U_{7/8}(du)=\int_0^\infty g(v^2)u(v)\,dv$.

Route (a), direct. $u(v)=\pi^{-1}\Re(\xi'/\xi)(\tfrac78+iv)$, with $\xi'/\xi$ from its gamma factor and $\zeta'/\zeta$ (Section~\ref{cx:sec:details}), needs $\zeta$ and $\zeta'$ inside the critical strip. The formula defines a signed density unconditionally; the announced half plane is what makes it the Thorin measure of \eqref{cx:eq:F78}. Its minimum on $[0,400]$ is $0.0055>0$ and its mass on $[0,400]$ is $201.63$ against $202$ zeros, the deficit being the Cauchy tails of width $3/8$ at the ends.

Route (b), deconvolution. Run the accelerated EM of Section~\ref{cx:sec:accel} on the level one data with the Cauchy kernel of width $1-\tfrac78=\tfrac18$. The condition number $e^{\pi/(8h)}$ is $e^{7.9}$ at $h=0.05$, mild, and $400$ steps reach a residual of $2\times10^{-10}$ (the author's run). Away from the edges of the data window it agrees with route (a) to $6\times10^{-4}$ absolutely and $0.33$ per cent relatively.

Route (c), zeros. Under the Riemann hypothesis, which is verified far beyond height $400$ \cite{PT21},
\[
u(v)=\sum_\rho\frac{(3/8)/\pi}{(3/8)^2+(v-\gamma_\rho)^2},
\]
summed over all zeros with both signs of $\gamma$. With the first $1000$ zeros and a tail from the zero density $\frac1{2\pi}\log\frac v{2\pi}$ above the last one, this agrees with route (a) to $9\times10^{-5}$ absolutely and $1.6$ per cent relatively, the relative discrepancy being the crude tail at the top of the range.

Finally the identity \eqref{cx:eq:F78} itself, with route (a) on the right. Minus the logarithm of the left side is $\log\xi(\tfrac78+\sqrt s)-\log\xi(\tfrac78)$ and the right side is $\int_0^\infty\log(1+s/v^2)\,u(v)\,dv$. At $s=1$ the left side is $0.040368$, the integral over $[0,400]$ is $0.038361$ and the tail from the zero density $0.002051$, total $0.040412$; at $s=4$ the three numbers are $0.126495$, $0.118335$, $0.008202$, total $0.126537$. The agreement is to the accuracy of the tail.

\subsection{What this gives the programme}

Route (b) produces the Thorin measure at level $7/8$ from data on the line $\Re s=1$, where the Euler product is in force, by a positivity preserving iteration with no reference to zeros and no continuation of $\zeta$ into the strip. The same computation runs at any $\theta'\in(1/2,7/8)$; the half plane adds at $7/8$ the guarantee that the limit is a Thorin measure rather than an information projection with a residual. Watching the residual at $\theta'<7/8$ is therefore the computational form of the question whether the zero free half plane extends, a diagnostic and not a certificate, since Proposition~\ref{cx:prop:what}(b) is qualitative.

\section{Pictures of the Thorin measures and the decision of positivity}\label{cx:sec:pictures}

\subsection{Pictures}

The densities $u_\theta$ on $[0,60]$ at $\theta=0.55,0.6,0.75,7/8,1$, each computed directly on its own line, show the Cauchy semigroup (Proposition~\ref{prop:struct}) from the side: as $\theta$ decreases to $\tfrac12$ the peaks at the zeros sharpen like $1/(\pi(\theta-\tfrac12))$ and the troughs between them deepen, and at $\theta=\tfrac12$ the picture becomes the zero measure.

The level $7/8$ measure on $[0,400]$ oscillates about $(2\pi)^{-1}\log(v/2\pi)$, the density of the zeros, as the average of a sum of Cauchy bumps must. In the variable $u=v^2$ the Thorin measure $U_{7/8}$ has density $g(u)=u_{7/8}(\sqrt u)/(2\sqrt u)$, and its two asymptotes make the Thorin conditions visible: $g(u)\sim u_{7/8}(0)/(2\sqrt u)$ as $u\to0$, so $\int_0^1|\log u|\,U_{7/8}(du)<\infty$, and $g(u)\approx\log(\sqrt u/2\pi)/(4\pi\sqrt u)$ on average as $u\to\infty$, so $\int_1^\infty u^{-1}U_{7/8}(du)<\infty$. The total mass is infinite, consistent with $F_{7/8}(s)=\xi(7/8)/\xi(7/8+\sqrt s)$ decaying faster than any power of $s$, in fact like $\exp(-\tfrac12\sqrt s\log\sqrt s)$.

\subsection{What positivity means}

\begin{theorem}[positivity]\label{cx:thm:positivity}
For $\tfrac12\le\theta<1$ the following are equivalent.
\begin{enumerate}
\item[(i)] $U_\theta$ is a positive Thorin measure, that is, $F_\theta(s)=\xi(\theta)/\xi(\theta+\sqrt s)$ is the Laplace transform of a generalised gamma convolution.
\item[(ii)] $\zeta$ has no zeros in $\Re s>\theta$.
\item[(iii)] $\Re(\xi'/\xi)(s)>0$ for $\Re s>\theta$; that is, $u_\sigma(v)>0$ for every $\sigma>\theta$ and every $v$.
\end{enumerate}
For $\theta=\tfrac12$ the equivalence of (ii) and (iii) is Lagarias's theorem \cite{Lagarias}.
\end{theorem}

\begin{proof}
(i) and (ii) are equivalent by Theorem~\ref{lv:thm:thorin}. (ii) implies (iii) term by term, since $\Re(\xi'/\xi)(s)=\sum_\rho(\sigma-\sigma_\rho)/|s-\rho|^2$, the sum taken over the zeros in pairs $\rho$, $1-\bar\rho$, and every $\sigma-\sigma_\rho$ is positive. For (iii) implies (ii), let $\rho_0=\sigma_0+i\gamma_0$ be a zero of multiplicity $m$ with $\sigma_0>\theta$. Near $\rho_0$, $\xi'/\xi(s)=m/(s-\rho_0)+h(s)$ with $h$ analytic, so $\Re(\xi'/\xi)(\sigma_0-\varepsilon+i\gamma_0)=-m/\varepsilon+O(1)$, which is negative for small $\varepsilon$, and $\sigma_0-\varepsilon>\theta$.
\end{proof}

Positivity on the single line $\Re s=\theta$, $u_\theta\ge0$, is necessary for (ii) but not sufficient. An off line zero $\rho_0$ puts a negative bump of height $1/(\pi(\sigma_0-\theta))$ at $v=\gamma_0$, and its mirror $1-\bar\rho_0$ a positive bump of height $1/(\pi(\theta+\sigma_0-1))$. The remaining zeros add a background which, once $\theta-\tfrac12$ exceeds the mean spacing of the zeros near $\gamma_0$, is about the zero density
\[
d(\gamma_0)=\frac1{2\pi}\log\frac{\gamma_0}{2\pi}.
\]
Heuristically, since the background grows with the height, a zero high enough up is invisible to pointwise positivity on a line to its left (an open question, Remark~\ref{rem:ch2:oneline}); at height $60$ the opposite holds (last column of Table~\ref{cx:tab:deficit}).

\subsection{The mass deficit}

The invariant that does not mask is the mass.

\begin{proposition}[mass deficit]\label{cx:prop:deficit}
Let $\tfrac12<\theta<1$ and $T>0$, and suppose no zero lies on the segment $\Re s=\theta$, $0\le\Im s\le T$, nor on the segment $\Im s=T$, $\theta\le\Re s\le1$. Let $M_\sigma(T)=\int_0^Tu_\sigma(v)\,dv$ and let $N(\theta,T)$ be the number of zeros with $\Re\rho>\theta$ and $0<\Im\rho<T$. Then
\begin{equation}\label{cx:eq:deficit}
2N(\theta,T)=M_1(T)-M_\theta(T)+\frac1\pi\Delta_{\mathrm{top}}(\theta,T),
\end{equation}
where $\Delta_{\mathrm{top}}(\theta,T)=-\int_\theta^1\Im(\xi'/\xi)(\sigma+iT)\,d\sigma$ is the change of $\arg\xi$ along the top side.
\end{proposition}

\begin{proof}
Apply the argument principle to $\xi$ on the rectangle with corners $\theta$, $1$, $1+iT$, $\theta+iT$, traversed counterclockwise. There are no zeros on $\Re s=1$, and $\xi$ is real and positive on $[\theta,1]$, so the bottom side contributes nothing. By the Cauchy Riemann equations $\partial_v\arg\xi(\sigma+iv)=\Re(\xi'/\xi)(\sigma+iv)=\pi u_\sigma(v)$, so the right side contributes $\pi M_1(T)$ and the left side, traversed downward, $-\pi M_\theta(T)$. The top side contributes $\Delta_{\mathrm{top}}$, since $\partial_\sigma\arg\xi=\Im(\xi'/\xi)$. The total is $2\pi N(\theta,T)$.
\end{proof}

The identity is the argument principle of Backlund and Turing; the Thorin picture supplies the bookkeeping, not a cheaper computation. As $\theta$ decreases to $\tfrac12$ with no zeros off the line below $T$, $M_\theta(T)$ tends to the number of zeros on the critical line below $T$, and \eqref{cx:eq:deficit} becomes the classical formula $N(T)=M_1(T)+\pi^{-1}\Delta_{\mathrm{top}}(\tfrac12,T)$, which is Backlund's count and holds with $N(T)$ the number of all zeros below $T$, with no hypothesis, when $T$ is not an ordinate.

Table~\ref{cx:tab:deficit} evaluates \eqref{cx:eq:deficit} at $T=100$ for the real $\xi$, and for a synthetic $\tilde\xi$ in which the consecutive zeros $\tfrac12+59.347i$ and $\tfrac12+60.832i$ and their conjugates are replaced by the off line quadruple $0.85\pm60.089i$, $0.15\pm60.089i$. The factor is a ratio of products of two quadratics, with the exact zeros in the denominator, symmetric under $s\mapsto1-s$ and positive on the real axis, so $\tilde\xi$ is entire of order one with the functional equation of $\xi$. Arguments were tracked on a grid of spacing $0.02$ in $v$ and $0.0005$ in $\sigma$ (the author's computation, not rerun).

\begin{table}[h]
\centering
\small
\begin{tabular}{cccccc}
\toprule
$\theta$ & $M_\theta(100)$ & $\Delta_{\mathrm{top}}/\pi$ & $2N$ for $\xi$ & $2N$ for $\tilde\xi$ & $\min\tilde u_\theta$\\
\midrule
0.55 & 29.0109 & $-0.1002$ & 0.0000 & 2.0000 & $-0.26$\\
0.60 & 29.0218 & $-0.0893$ & 0.0000 & 2.0000 & $-0.55$\\
0.65 & 29.0328 & $-0.0784$ & 0.0000 & 2.0000 & $-0.94$\\
0.70 & 29.0438 & $-0.0674$ & 0.0000 & 2.0000 & $-1.51$\\
0.75 & 29.0548 & $-0.0563$ & 0.0000 & 2.0000 & $-2.57$\\
0.80 & 29.0659 & $-0.0452$ & 0.0000 & 2.0000 & $-5.38$\\
0.875 & 29.0827 & $-0.0284$ & 0.0000 & 0.0000 & $0.006$\\
0.90 & 29.0884 & $-0.0228$ & 0.0000 & 0.0000 & $0.006$\\
0.95 & 29.0997 & $-0.0114$ & 0.0000 & 0.0000 & $0.007$\\
1.00 & 29.1112 & $0$ & 0.0000 & 0.0000 & $0.007$\\
\bottomrule
\end{tabular}
\caption{The mass deficit at $T=100$, where $M_1(100)=29.1112$ and $N(100)=29$ (the author's computation, not rerun).}\label{cx:tab:deficit}
\end{table}

For the real $\xi$ the deficit vanishes to four decimals at every level of the table, and $M_\theta(100)$ decreases almost linearly, with slope $0.223$ per unit of $\theta$, towards $29.000$ as $\theta\to\tfrac12$: the smoothed zero count becomes the zero count. For $\tilde\xi$ the deficit is $2$ exactly when $\theta<0.85$ and $0$ above, as the table shows; at this height the synthetic density is also visibly negative near $60$ at every level below $0.85$.

\begin{corollary}[the level mass function]\label{cx:cor:levelmass}
Let $T>0$ not be the ordinate of a zero. For $\theta\in(\tfrac12,1]$ not the real part of a zero below height $T$,
\[
M_\theta(T)=M_1(T)+\frac1\pi\Delta_{\mathrm{top}}(\theta,T)-2N(\theta,T).
\]
The first two terms form a real analytic function of $\theta$, since $\xi$ has no zero on the line $\Im s=T$; its derivative is
\[
\frac1\pi\Im\frac{\xi'}{\xi}(\theta+iT).
\]
The last term is a step function which drops by $2m$ as $\theta$ decreases through $\Re\rho$, for each zero $\rho$ of multiplicity $m$ with $\Re\rho>\tfrac12$ and $0<\Im\rho<T$. Hence $\theta\mapsto M_\theta(T)$ extends continuously to $(\tfrac12,1]$ if and only if every zero below height $T$ lies on the critical line, and its jumps occur exactly at the real parts of the off line zeros, of size twice the total multiplicity of the zeros with that real part and ordinate below $T$.
\end{corollary}

At $T=100$, for $\xi$ the masses on the ten vertical lines of the table agree with the smooth part, computed from one horizontal sweep, to $2.3\times10^{-5}$; for $\tilde\xi$ they lie exactly $2.0000$ below it for $\theta<0.85$ and on it above (the author's computation, not rerun). The corollary does not locate real parts more cheaply than the argument principle, since the jumps are seen only by computing along vertical lines.

\subsection{Positivity of $U_{7/8}$}

Given the zero free half plane $\Re s>7/8$ (announced, not refereed, as above), Theorem~\ref{cx:thm:positivity} makes $U_{7/8}$ positive: every term of $u_{7/8}=\sum_\rho P_{7/8-\sigma_\rho}(v-\gamma_\rho)$, with $P_a$ the Cauchy density of width $a$ and $P_0=\delta_0$, is nonnegative, a zero on the line $\Re s=7/8$, if there were one, contributing an atom. The computation displays the margin. On $[0,400]$ the minimum of $u_{7/8}$ is $0.00551$, attained at $v=0$, where $u_{7/8}(0)=\pi^{-1}(\xi'/\xi)(7/8)=0.005514663528$; this value is positive unconditionally, because $\xi(\tfrac12+t)=\int\Phi(x)\cosh(tx)\,dx$ with Riemann's kernel $\Phi$ positive, so $\xi$ increases on $(\tfrac12,\infty)$. Away from the origin the minimum on $[15,400]$ is $0.0272$, at $v=17.5$ between the first two zeros; with the second difference estimate $\max|u''|\approx11.8$ on the grid of spacing $0.05$, the grid minimum less $h^2/8\cdot\max|u''|$ is $0.0235$. Theorem~\ref{cx:thm:certified} certifies the positivity of $u_{7/8}$ in ball arithmetic on $[0,1420]$, and Theorem~\ref{cx:thm:globalmin} certifies the lower bound $0.0272$ for $u_{7/8}$ on $[13.5,1420]$.

\subsection{The measure $U_{3/4}$ below level 7/8}

At every height the positivity of $U_{3/4}$ is, by Theorem~\ref{cx:thm:positivity}, the open zero free half plane $\Re s>3/4$. We computed it as at $7/8$, the direct route on a grid of spacing $0.02$ ($20{,}001$ points) since the peaks have width $1/4$. The table is the author's computation, not rerun apart from the value at the origin.

\begin{center}
\small
\begin{tabular}{ll}
\toprule
quantity & value\\
\midrule
$\min u_{3/4}$ on $[0,400]$ & $0.00368$ at $v=0$\\
$\min u_{3/4}$ on $[15,400]$ & $0.0182$ at $v=17.5$\\
$\max|u''|$; margin & $40.5$; \ $0.0162$\\
mass on $[0,400]$ & $201.60$ against $N(400)=202$\\
zeros plus tail minus direct & $5.7\times10^{-5}$\\
EM: residual; relative error & $1.5\times10^{-9}$; \ $1.1$ per cent\\
EM density: minimum & $4.0\times10^{-3}$\\
$2N(\tfrac34,400.9)$ & $202.1778-202.0891-0.0887=0.0000$\\
\bottomrule
\end{tabular}
\end{center}

\noindent The EM run used the Cauchy kernel of width $1/4$, condition number $e^{\pi/(4h)}=e^{15.7}$ at $h=0.05$, and $1500$ accelerated steps; the height $T=400.9$ for the deficit lies between the zeros $399.985$ and $401.839$. The minimum at the origin equals $\pi^{-1}(\xi'/\xi)(3/4)=0.003676904191$. So $u_{3/4}$ is positive on $[0,400]$ and its mass deficit to height $400.9$ is zero: no zero lies to the right of $3/4$ below that height, as was known far higher \cite{PT21}; Theorem~\ref{cx:thm:certified} certifies this on $[0,1420]$.

\subsection{All the way down: $U_{1/2}$}\label{cx:sec:half}

At $\theta=\tfrac12$ the densities collapse. Since $\xi(\tfrac12+iv)=\Xi(v)$ is real, $\Re(\xi'/\xi)(\tfrac12+iv)=0$ away from the zeros, and the measure is purely atomic: $\pi\,\mathfrak U_{1/2}=\partial_v\arg\xi(\tfrac12+iv)$, which jumps by $\pi m$ at a zero of multiplicity $m$ on the line. Under the Riemann hypothesis $\mathfrak U_{1/2}=\sum_\rho\delta_{\gamma_\rho}$ and $U_{1/2}=\sum_{\gamma>0}m_\gamma\delta_{\gamma^2}$, the Thorin measure of the reciprocal $\xi$ function around which the book is organised. It can be computed three ways; the author's last figure shows the cumulative masses $M_\theta(v)$ at $\theta=1,0.75,0.55$ converging to the zero staircase $N(v)$, and the EM backward flow at $\theta'=\tfrac12$ on a logarithmic scale with the sign changes of $Z$ marked.

Directly at level $\tfrac12$, the atoms of odd multiplicity are the sign changes of Hardy's function $Z(v)=e^{i\vartheta(v)}\zeta(\tfrac12+iv)$, which is real with the sign of $\Xi(v)$ up to a fixed sign. A scan of $[1,400.5]$ at spacing $0.02$, followed by root finding, gives $202$ atoms below height $400$, agreeing with the tabulated zeros to machine precision, the classical computation, included for comparison.

By the backward flow at $\theta'=\tfrac12$ from level one, the runs of Section~\ref{cx:sec:accel} apply: on $[-400,400]$ with $3000$ accelerated steps the peaks lie within $0.010$ of the zeros, $0.0055$ on average, with basin masses $1.0013\pm0.0313$ ($\pm0.0061$ for the interior basins), and to height $1415$ the basin masses are $1.0007\pm0.0053$, the unresolved pair at $1329.125$ carrying mass $2.001$. The author's centroid statements (the basin centroids lie on average $0.0014$ and at most $0.081$ from the zeros, the density between the atoms falls to $10^{-12}$) were not rerun.

Is $U_{1/2}$ positive? That is the Riemann hypothesis (Theorem~\ref{lv:thm:thorin}). To certify positivity up to height $T$, the sign changes, which count zeros of odd multiplicity on the line, need the total count $N(T)$, which Backlund's count supplies from level one data:
\[
N(T)=M_1(T)+\frac1\pi\Delta_{\mathrm{top}}(\tfrac12,T).
\]
At $T=400.9$ this gives $202.1778-0.1778=202.0000$, and there are $202$ sign changes of $Z$ below $400.9$. Every zero below height $400.9$ is therefore on the line and simple, which Theorem~\ref{cx:thm:certified} certifies to height $1419.9$, and $\mathfrak U_{1/2}$ restricted to $|v|<400.9$ is the sum of unit atoms at the ordinates. This is Turing's verification \cite{Turing53} in Thorin bookkeeping.

\subsection{Can positivity be decided?}\label{cx:sec:decide}

By level. For $\theta\ge1$ positivity is an unconditional theorem, and for $\theta\ge7/8$ it holds conditionally on the announced half plane; the computation only displays it.

For $\tfrac12<\theta<7/8$ and finite $T$, positivity up to height $T$, $N(\theta,T)=0$, is decidable: compute the deficit \eqref{cx:eq:deficit} to accuracy better than one half in interval arithmetic, provided no zero lies on the boundary of the rectangle, which is known within the verified range. For $\theta=\tfrac12$ Backlund's count and the sign changes decide it. Theorem~\ref{cx:thm:certified} does this in ball arithmetic at $\theta=3/4$ and $T=1419.9$.

For $\tfrac12\le\theta<7/8$ and all heights, positivity is the zero free half plane $\Re s>\theta$, and for $\theta=\tfrac12$ the Riemann hypothesis. If false it is refutable by a finite computation: an offending zero lies in a finite rectangle, over which the deficit is a positive even integer. If true it is not confirmable by finitely many heights without an argument for the rest. That is why the half plane at $7/8$ is claimed by a proof, and why no computation in this book settles the positivity of $U_\theta$ below $7/8$.

\subsection{Computational details}\label{cx:sec:details}

The author's computations used Python with \texttt{mpmath} \cite{mpmath} at 15 significant digits for $\zeta$, $\zeta'$, $\log\Gamma$, $\psi$ and Hardy's $Z$, \texttt{numpy} for the fast Fourier transforms and the EM iterations. The first $1000$ zeros were taken from \texttt{mpmath.zetazero}. The recomputations of this book used Arb through \texttt{python-flint} \cite{Arb,pythonflint} at 96 to 160 bits.

Densities. $u_\theta(v)=\pi^{-1}\Re(\xi'/\xi)(\theta+iv)$ was evaluated from
\[
\frac{\xi'}{\xi}(s)=\frac1s+\frac1{s-1}-\frac{\log\pi}2+\frac12\psi\Bigl(\frac s2\Bigr)+\frac{\zeta'}{\zeta}(s).
\]
The grid spacing was chosen so that each Cauchy peak of width $\theta-\tfrac12$ is covered by at least ten points: $0.05$ at $\theta=7/8$ and $1$, $0.02$ at $\theta=3/4$, $0.01$ at $\theta=0.55$ and $0.6$.

Arguments. The masses $M_\theta(T)$ and the top correction were obtained by continuous variation of $\Im\log\xi$, computed as $\Im[\log s+\log(s-1)-\tfrac s2\log\pi+\log\Gamma(s/2)+\log\zeta(s)]$ and unwrapped along the path. The unwrapping is valid when the phase increment per step is below $\pi$; since $\partial_v\arg\xi=\pi u_\theta$ and $\pi u_\theta\le1/(\theta-\tfrac12)$ plus a background of order $\log v$, a step of $0.02$ gives increments below $0.5$ radians at $\theta=0.55$. The integer deficits came out within $10^{-4}$ of integers in every case.

Deconvolution. The matrix is \eqref{cx:eq:Amatrix} with cell densities as unknowns, the data grid $[-400,400]$ by evenness of $u_1$, the unknown grid $[-440,440]$, Toeplitz products by FFT, and the conventions for residuals (with the factor $h$; without it every value is $1/h$ times larger), steps, peaks and basins are those of Section~\ref{cx:sec:accel}.

The author's timings on one core, with \texttt{mpmath} for the data, orders of magnitude only (with Arb the data layer does not dominate; Section~\ref{cx:sec:thousand}):

\begin{center}
\small
\begin{tabular}{lr}
\toprule
computation & seconds\\
\midrule
$u_1$ on $[0,400]$, $8001$ points & 103\\
$u_{7/8}$ on $[0,400]$, $8001$ points & 100\\
$u_{3/4}$ on $[0,400]$, $20{,}001$ points & 239\\
$u_\theta$ on $[0,60]$ at five levels & 60\\
arguments on ten vertical lines to $T=100$ and one sweep & 148\\
arguments to $T=400.9$ at $\theta=3/4$ and $1$, two sweeps & 207\\
EM at width $1/8$, $400$ steps, $17{,}601$ unknowns & 3\\
EM at width $1/4$, $1500$ steps & 11\\
sign changes of $Z$ on $[1,400.5]$ & 122\\
$u_1$ on $[0,1450]$, $58{,}001$ points; EM to height $1450$ & 1375; 151\\
\bottomrule
\end{tabular}
\end{center}

The positivity of $u_{3/4}$ and $u_{7/8}$ on $[0,1420]$ and the zero counts above are certified in ball arithmetic in Section~\ref{cx:sec:certify}; the other numbers above are floating point.

\subsection{Certification in ball arithmetic}\label{cx:sec:certify}

The positivity statements and integer counts above were obtained in floating point. Here they are repeated in the ball arithmetic of Arb \cite{Arb}, where each quantity is an interval guaranteed to contain the true value, including all rounding and truncation errors. Two lemmas reduce a line or segment to finitely many ball evaluations.

\begin{lemma}[cells]\label{cx:lem:cell}
Let $L=\xi'/\xi$, $u(v)=\pi^{-1}\Re L(\theta+iv)$, and let $\xi$ have no zero in the closed cell $\{\theta+iv:|v-m|\le r\}$. Then for $|t|\le r$
\[
u(m+t)\ \ge\ u(m)-r\,|u'(m)|-\frac{r^2}{2\pi}\sup_{|v-m|\le r}|L''(\theta+iv)|,
\]
where $u'(m)=-\pi^{-1}\Im L'(\theta+im)$.
\end{lemma}

\begin{proof}
Taylor's theorem with the Lagrange remainder, using $\partial_vL(\theta+iv)=iL'$ and $\partial_v^2L(\theta+iv)=-L''$, so that $|u''|\le|L''|/\pi$.
\end{proof}

In practice $u(m)$ and $u'(m)$ are evaluated at the exact midpoint, where the balls are tight, and the supremum is bounded by evaluating the power series of $L$ on the whole cell as a ball, which overestimates but is multiplied by $r^2/2$. A cell whose bound is not certifiably positive, or on which the ball enclosure of $\xi$ contains $0$, is bisected.

\begin{lemma}[segment]\label{cx:lem:segment}
Let $[a,b]$ be a segment of length $d$ with
\[
d\,\sup_{s\in[a,b]}|\xi'(s)|<|\xi(a)|.
\]
Then $\xi$ has no zero on the segment, and the continuous change of $\arg\xi$ along it is the principal value $\mathrm{Arg}(\xi(b)/\xi(a))$, which lies in $(-\pi/2,\pi/2)$.
\end{lemma}

\begin{proof}
For $s\in[a,b]$, $|\xi(s)-\xi(a)|\le d\sup|\xi'|<|\xi(a)|$, so $\xi(s)/\xi(a)$ lies in the open disk of radius less than one about $1$, which lies in the half plane $\Re w>0$. The continuous argument of $\xi(s)/\xi(a)$ therefore stays in $(-\pi/2,\pi/2)$ along the segment and ends at the principal value.
\end{proof}

The cells have exact dyadic endpoints, the segments exact dyadic endpoints except for the corner $T^*=14199/10$, which is one ball shared by all paths, so that the contours close and every enclosure contains the value at the exact corner.

\begin{theorem}[certified]\label{cx:thm:certified}
Let $T^*=1419.9$. Then:
\begin{enumerate}
\item[(i)] $u_{3/4}(v)>0$ and $u_{7/8}(v)>0$ for $0\le v\le 1420$;
\item[(ii)] $N(T^*)=1000$;
\item[(iii)] $N(\tfrac34,T^*)=0$, with $M_{3/4}(T^*)=1000.01952570266\ldots$ to within $10^{-20}$;
\item[(iv)] $\Xi$ has $1000$ certified sign changes in $(0,T^*)$.
\end{enumerate}
Consequently every zero of $\zeta$ with $0<\Im\rho<T^*$ lies on the critical line and is simple; the measure $\pi^{-1}\partial_v\arg\xi(\tfrac12+iv)$ of Section~\ref{cx:sec:half} restricted to $0\le v<T^*$, which under the Riemann hypothesis is the Thorin measure $\mathfrak U_{1/2}$, is the sum of $1000$ unit atoms at the ordinates; and $u_{3/4}$ and $u_{7/8}$ have positive densities there, by (i).
\end{theorem}

\begin{proof}
{\emergencystretch=8em Computed, computer assisted, in the ball arithmetic of Arb at 96 bits and again at 160 bits; two standard facts are used and not certified, that $\xi>0$ on $(\tfrac12,\infty)$ and Backlund's relation $\pi N(T)=\Delta\arg\xi$ along $2\to2+iT\to\tfrac12+iT$. For (i), Lemma~\ref{cx:lem:cell} on the dyadic cells of width $2^{-7}$ over $[0,1420]$, bisected once where necessary, gives $182{,}409$ cells for $\theta=3/4$ and $181{,}781$ for $\theta=7/8$ at 96 bits ($182{,}512$ and $181{,}783$ at 160 bits; the counts belong to the method and the grid, not to $u$); every certified lower bound is positive, and in both runs and at both levels it is at least $4\times10^{-5}$. The least bound of the certificate depends on the grid and the precision and is not the minimum of $u$, which for $u_{3/4}$ and $u_{7/8}$ on $[0,1420]$ is attained near $0$, where $u_{3/4}(0)=\pi^{-1}(\xi'/\xi)(3/4)=0.003676904191347$ and $u_{7/8}(0)=0.005514663527790$ (balls of radius below $10^{-17}$). For (ii), Lemma~\ref{cx:lem:segment} along the path $2\to2+iT^*\to\tfrac12+iT^*$, which gives $\pi N(T^*)$ as the change of argument because $\xi$ is real and positive on $[\tfrac12,2]$ and real on the critical line: the enclosure contains $1000$ and has width below $10^{-20}$ ($11{,}470$ segments on $\Re s=2$ and $96$ pieces of length $1/64$ on the top side at 96 bits). For (iii), Lemma~\ref{cx:lem:segment} around the rectangle $[\tfrac34,2]\times[0,T^*]$, using Proposition~\ref{cx:prop:deficit} with the right side moved from $\Re s=1$ to $\Re s=2$, which changes nothing since there are no zeros with $\Re s\ge1$ and keeps the balls away from the pole of $\zeta$: the enclosure of $2N(\tfrac34,T^*)$ has radius below $10^{-20}$ about $0$ and contains the single integer $0$ ($32{,}387$ segments on $\Re s=\tfrac34$), and $M_{3/4}(T^*)=1000.0195257026609355$ with the same twenty digits at both precisions. For (iv), the signs of $\Xi(v)=\xi(\tfrac12+iv)$ at the $70{,}996$ points $v=k/50$, $k=0,\dots,70995$, are certified, none undetermined, with $\Xi(0)>0$, and change $1000$ times. Since the sign changes give at least $1000$ distinct zeros of odd multiplicity on the line and (ii) gives exactly $1000$ zeros in all, each zero is on the line and simple. The statement about $\pi^{-1}\partial_v\arg\xi$ is (ii) with (iv); the positivity statement is (i); (iii) is implied by (ii) and (iv) and is an independent check.\par}
\end{proof}

\begin{table}[h]
\centering
\small
\begin{tabular}{lll}
\toprule
certificate & enclosure at 96 bits (160 bits) & size\\
\midrule
(i) $u_{3/4}>0$ on $[0,1420]$ & least bound at least $4\times10^{-5}$ & $182{,}409$ cells\\
(i) $u_{7/8}>0$ on $[0,1420]$ & least bound at least $4\times10^{-5}$ & $181{,}781$ cells\\
(ii) $N(1419.9)$ & $1000\pm4\times10^{-22}$ ($\pm2\times10^{-41}$) & $11{,}470+96$ segments\\
(iii) $2N(\tfrac34,1419.9)$ & $0\pm1.4\times10^{-21}$ ($\pm7\times10^{-41}$) & $32{,}387$ segments\\
(iv) signs of $\Xi$ & $1000$ changes, none undetermined & $70{,}996$ points\\
\bottomrule
\end{tabular}
\caption{The certificates of Theorem~\ref{cx:thm:certified}, as recomputed.}\label{cx:tab:certified}
\end{table}

\noindent The author's run to height $400.9$ certifies $N(400.9)=202$, an enclosure of radius below $10^{-16}$ for $2N(\tfrac34,400.9)$, $202$ sign changes, and both positivity statements on $[0,400]$, in under two minutes (not rerun).

Parts (ii) to (iv) are known \cite{PT21}; the theorem adds part (i), the positivity of the Thorin densities at levels $3/4$ and $7/8$ on $[0,1420]$ as a computer assisted statement, fast enough for a laptop.

\section{The limits of the computation}\label{cx:sec:limits}

Can these computations be pushed past what is known? For $\zeta$, no.

Height. Theorem~\ref{cx:thm:certified} reaches height $1419.9$ in minutes. Platt and Trudgian verified the Riemann hypothesis rigorously to height $3\times10^{12}$ \cite{PT21}, about $1.2\times10^{13}$ zeros, combining Turing's method with the fast multi evaluation of Odlyzko and Sch\"onhage \cite{OS88}, a few operations per zero. The mass deficit is Turing's method in Thorin masses, so its known cost is the same, of order $T^{1+\varepsilon}$; matching the record would mean rebuilding that machinery.

All heights. Below $7/8$, positivity of $U_\theta$ at every height is the zero free half plane $\Re s>\theta$, which a computation stopping at a finite height $H$ cannot confirm (Section~\ref{cx:sec:decide}).

Combining the two gains nothing either. Below $3\times10^{12}-\tfrac12$ every zero is on the line, so $U_\theta$ restricted to that range is known to be positive for every $\theta>\tfrac12$ (Corollary~\ref{cor:dbeta:window}). Above it, the announced half plane gives positivity for $\theta\ge7/8$, conditionally, and nothing below. Every computation in these chapters lies inside what is already known. The frontier that moves the half plane is analytic: zero density estimates such as those of Guth and Maynard \cite{GM26}.

What computation can still add is narrower: two results about the Thorin densities themselves, one unconditional (Proposition~\ref{bp:prop:v1min}) and one conditional on the announced half plane, whose Riemann hypothesis form is Theorem~\ref{thm:dbeta:minall}(vi). They rest on a lemma.

\begin{lemma}[pairing]\label{cx:lem:pairing}
Let $\nu=\sum_\rho\delta_{\gamma_\rho}$ be the counting measure of the ordinates of the nontrivial zeros, with multiplicity and both signs, and $P_a$ the Cauchy density of width $a$, with $P_0=\delta_0$. If $\zeta$ has no zeros in $\Re s>\theta_0$, where $\tfrac12\le\theta_0\le1$, then for every $\theta\ge\theta_0$, as measures,
\[
\mathfrak U_\theta\ \ge\ \tfrac12\,P_{2\theta-1}*\nu .
\]
In particular the bound holds unconditionally for $\theta\ge1$, and for $\theta\ge7/8$ it holds conditionally on the announced half plane. Under the Riemann hypothesis $\mathfrak U_\theta=P_{\theta-1/2}*\nu$ exactly.
\end{lemma}

\begin{proof}
Group the zeros into orbits under $\rho\mapsto1-\bar\rho$, which preserves the ordinate. An orbit of two zeros $\sigma+i\gamma$, $1-\sigma+i\gamma$ contributes $P_a+P_b$ at the point $\gamma$, with widths $a=\theta-\sigma$ and $b=\theta-1+\sigma$, both nonnegative because $\theta\ge\theta_0\ge\max(\sigma,1-\sigma)$, and $a+b=2\theta-1$. For fixed $x\ne0$ write $\pi P_a(x)=a\,h(a)$ with $h(a)=1/(a^2+x^2)$, which is decreasing in $a\ge0$. Then
\[
a\,h(a)+b\,h(b)\ \ge\ a\,h(a+b)+b\,h(a+b)=(a+b)\,h(a+b),
\]
that is $P_a(x)+P_b(x)\ge P_{a+b}(x)=\tfrac12\cdot2P_{2\theta-1}(x)$, which is the claim for the two zeros of the orbit (the two zeros count twice in $\nu$); at $x=0$ the inequality holds as measures because a width zero term is an atom. An orbit of one zero, on the critical line, contributes $P_{\theta-1/2}$, and $P_{\theta-1/2}(x)\ge\tfrac12P_{2\theta-1}(x)$ by the same monotonicity with $a=\theta-\tfrac12$, $a+b=2a$. Summing over orbits gives the lemma.
\end{proof}

With the cell certificate and an explicit zero count, the lemma proves the following.

\begin{theorem}[global minimum]\label{cx:thm:globalmin}
(i) Unconditionally, for every real $v\ne0$,
\[
u_1(v)\ >\ u_1(0)=\frac1\pi\Bigl(1+\frac{\gamma_E}2-\frac12\log4\pi\Bigr)=0.0073515924923\ldots,
\]
where $\gamma_E$ is Euler's constant; this is Proposition~\ref{bp:prop:v1min}.
(ii) If $\zeta$ has no zeros in $\Re s>7/8$ (the announced half plane), then, as measures, $\mathfrak U_{7/8}\ge u_{7/8}(0)\,dv$ with $u_{7/8}(0)=\pi^{-1}(\xi'/\xi)(7/8)=0.0055146635278\ldots$, and $u_{7/8}(v)>u_{7/8}(0)$ wherever $v\ne0$ is not the ordinate of a zero on the line $\Re s=7/8$.
\end{theorem}

\begin{proof}
Part (i) is Proposition~\ref{bp:prop:v1min}; the argument below is the same in three ranges and gives (ii) as well. It is computed, computer assisted. Write $\theta$ for $1$ or $7/8$; both functions are even. Three ranges.

$|v|\le13.5$. Here $u_\theta=\sum_\rho P_{\theta-\sigma_\rho}(v-\gamma_\rho)$ with every width in $[0,1)$ for $\theta=1$, and in $[0,\tfrac34]$ for $\theta=7/8$ under the hypothesis. The Cauchy density $P_a(x)$ is convex for $|x|\ge a/\sqrt3$, and every zero has $|\gamma_\rho|\ge14.1347$ (Theorem~\ref{cx:thm:certified}), so $|v-\gamma_\rho|\ge0.63>1/\sqrt3$. Hence $u_\theta$ is a convergent sum of convex functions, strictly convex and even on $[-13.5,13.5]$, with its unique minimum at $0$.

$13.5\le|v|\le1420$. Lemma~\ref{cx:lem:cell} with the threshold $u_\theta(0)$ in place of $0$: on $175{,}813$ cells of radius $0.004$ for each of $\theta=1$ and $\theta=7/8$, with no bisection, every certified lower bound exceeds the upper end of the ball for $u_\theta(0)$, the smallest being $0.0360$ at $v=17.49$ for $\theta=1$ and $0.0272$ at $v=17.49$ for $\theta=7/8$ (the margins over $u_\theta(0)$ are $0.0287$ and $0.0217$). The computation evaluates $\xi'/\xi$ on the line and uses no hypothesis.

$|v|\ge1420$. By Lemma~\ref{cx:lem:pairing}, unconditional at $\theta=1$ and conditional at $7/8$, $u_\theta(v)\ge\tfrac12\sum_\gamma P_{2\theta-1}(v-\gamma)\ge\tfrac12\,n(v)\,P_{2\theta-1}(8)$, where $n(v)$ is the number of ordinates in $(v-8,v+8]$ (positive ordinates, with multiplicity). Hasanalizade, Shen and Wong \cite{HSW22} (the form used in Proposition~\ref{bp:prop:v1min}) give $|N(T)-\tfrac T{2\pi}\log\tfrac T{2\pi e}-\tfrac78|\le E(T)$ with $E(T)=2.5167+1/(50T)$ for $e\le T\le X_0=30\,610\,046\,000$ and $E(T)=0.1095\log T+0.2042\log\log T+3.0305+1/(50T)$ for $T>X_0$. Hence $n(v)\ge D(v)-E(v+8)-E(v-8)$ with $D(v)=M(v+8)-M(v-8)$ and $M(T)=\tfrac T{2\pi}\log\tfrac T{2\pi e}$. The right side is $8.7678\ldots$ at $v=1418.9$, so more than $8.767$, and it increases with $v$ as long as $v+8\le X_0$, because $M'(T)=\log(T/2\pi)/(2\pi)$ increases, so $D$ increases, and $E$ decreases; for larger $v$ it is at least $44$ at $X_0-8$ and increases (a numerical check of the derivative, in the recomputation). So $n(v)\ge9$ for $v\ge1420$, since $n$ is an integer, and $u_\theta(v)\ge4.5\,P_{2\theta-1}(8)$, which is $0.0220$ for $\theta=1$ and $0.0166$ for $\theta=7/8$, more than twice $u_\theta(0)$ in both cases. The conclusion would survive a loss of five ordinates at $\theta=1$ and of six at $\theta=7/8$.
\end{proof}

The same argument, with the Riemann hypothesis in place of the zero free half plane, covers every level at once, and the conclusion turns out to be equivalent to the hypothesis.

\begin{theorem}[a Thorin form of the Riemann hypothesis]\label{cx:thm:rhmin}
The Riemann hypothesis holds if and only if, for every $\theta\in(\tfrac12,1]$, the Thorin density $u_\theta$ attains its global minimum at $v=0$; that is,
\[
\Re\frac{\xi'}{\xi}(\theta+it)\ \ge\ \frac{\xi'}{\xi}(\theta)\qquad\text{for all real }t\text{ and }\tfrac12<\theta\le1 .
\]
Under the hypothesis the minimum is attained only at $t=0$.
\end{theorem}

\begin{proof}
For $\theta\in(\tfrac12,1)$ the forward direction, with the strict inequality, is Theorem~\ref{thm:dbeta:minall}(vi), which under the hypothesis holds at every ordinate, and for $\theta=1$ it is Proposition~\ref{bp:prop:v1min}. Conversely, if a zero $\rho_0=\sigma_0+i\gamma_0$ has $\sigma_0>\tfrac12$, take $\theta=\sigma_0-\varepsilon$; then $\Re(\xi'/\xi)(\theta+i\gamma_0)=-m/\varepsilon+O(1)<0$, while $(\xi'/\xi)(\theta)>0$ unconditionally, because $\xi(\tfrac12+x)=\int\Phi(y)\cosh(xy)\,dy$ with Riemann's kernel $\Phi$ positive increases for $x>0$. So the minimum is not at $0$.
\end{proof}

Theorem~\ref{cx:thm:uniform} below extends Theorem~\ref{cx:thm:globalmin}(ii) to every level in $[7/8,1]$. Theorem~\ref{cx:thm:rhmin} sits beside the reformulations of Lagarias \cite{Lagarias}, $\Re(\xi'/\xi)>0$ on $\Re s>\tfrac12$, and of Sondow and Dumitrescu \cite{SD10}, that $|\xi(\sigma+it)|$ increases in $\sigma$ for $\sigma>\tfrac12$, and strengthens both under the hypothesis: $\partial_\sigma\log|\xi(\sigma+it)|$ is not merely positive but smallest on the real axis, at every level in $(\tfrac12,1]$. Its forward direction at $\theta=1$ is computer assisted (Proposition~\ref{bp:prop:v1min}), and it is a reformulation, not progress towards the hypothesis.

The certificates transfer to any completed $L$ function with a ball arithmetic evaluator (Dirichlet $L$ functions, including the Dirichlet beta function, that of Ramanujan's $\Delta$, Dedekind zeta functions); whether they would add anything to the published verifications is not examined here. The Thorin measures at levels $7/8$ and $3/4$ and their margins are tables of numbers, not theorems about zeros.
